\chapter{Conclusion}
\label{ch:conclusion}

\openingpara
{This dissertation advances affective Human--Building Interaction} {as an orientation for specifying and investigating the situated complexity of indoor experience and bringing this understanding into the design of future interactive indoor environments. It maps the experiential agenda of research on interactive indoor environments, examines how environmental conditions acquire significance through affect, develops multisensory biophilic design knowledge, and investigates one sound-based instantiation through an AI-transformed \sidenote{Throughout this dissertation, I use the term \emph{transformed} to refer to AI-transformed audio. The RAVE model receives an incoming audio signal and generates a corresponding output through the learned nature-sound representation.} office soundscape. The studies show that experiential complexity lies both in what occupants feel and in how they sense, interpret, anticipate, negotiate, and adjust to their surroundings over time. This chapter answers the dissertation's research questions, brings its contributions together, and develops their implications for HBI.}

\section{research questions and dissertation contributions}
\label{sec:returning-problem}

The dissertation addressed a problem identified within HBI: human experience is central to the study of interactive buildings \cite{alavi2019introduction,alavi2017comfort}, while the forms and relations encompassed by this experience require greater understanding and specification (cf.~\autoref{ch:study-one}). This section returns to the four research questions through which this problem was addressed, states the corresponding dissertation contributions, and then responds to the overarching research question.

\clearpage
\paragraph{\rqref{scoping}.}

\emph{What forms of human experience are addressed in research on interactive indoor environments, and how does this research seek to shape them through design and technology?}

The scoping review presented in \autoref{ch:study-one} identifies a broad
experiential agenda across research on interactive indoor environments.
Across 192 papers, it distinguishes 11 categories of human experience:
\textit{Aesthetic Experience}, \textit{Bodily Experience},
\textit{Comfort}, \textit{Democratic Experience}, \textit{Explorative
and/or Future and Speculative Experiences}, \textit{Emotional
Experience}, \textit{Experience of Relatedness}, \textit{Inclusion
Experience}, \textit{Learning and Teaching Experience}, \textit{Sense
of Orientation}, and \textit{Sense of Place-making}. This breadth is
distributed across different disciplinary perspectives, methodologies,
and design approaches.

The review further identifies 20 design mechanisms through which these
experiences are pursued. These include mechanisms such as
\textit{Spatial Design}, \textit{Community and Social Connection},
\textit{Participation and Co-design}, \textit{Somaesthetic Design},
\textit{Technological Immersion}, \textit{Personalisation}, and
\textit{Biophilic Design}. Some mechanisms support several forms of
experience, while others are more closely associated with particular
experiential aims
(cf.~\autoref{sec:ch2-findings-design-mechanisms}). The review thus
specifies both the forms of experience addressed across the field and
the different design approaches through which researchers have sought
to shape them.

Two typologies specify the technological means through which this work is undertaken. The first comprises seven intervention types: \textit{Interactive Displays}, \textit{Audio and Lighting Technologies}, \textit{Environmental Sensing}, \textit{Feedback Systems}, \textit{Augmented and Virtual Reality}, \textit{Physiological Sensing}, and \textit{Shape-Changing and Kinetic Interfaces}. The second positions technological interventions across the temporal layers of buildings by adapting \citet{brand1995buildings}'s shearing layers model and introducing \textit{Extended Reality Space} as an additional
layer. This temporal analysis shows that interactive technologies are
concentrated in the more adaptable parts of buildings, with comparatively little work engaging longer-lasting architectural layers (cf.~\autoref{sec:ch2-findings-technology}).

The first dissertation contribution is therefore a structured account of what forms of experience have been addressed, the design mechanisms through which
they have been pursued, and the technological interventions through
which these approaches have been realised. This mapping establishes the experiential landscape for the subsequent studies. It also identifies particular directions for further examination: \textit{Emotional Experience} constitutes a substantial strand of the reviewed literature, while \textit{Biophilic Design} appears as a less frequently represented design mechanism. The dissertation subsequently examines these directions in greater depth.

\clearpage  


\paragraph{\rqref{afi}.}

\emph{What does an affective lens reveal about how occupants experience
and interpret interactive indoor environments, and how can these
insights inform HBI research and design?}


The two occupant studies presented in \autoref{ch:study-two} adopt
Affective Interaction (AfI) from HCI as the dissertation's affective
lens.\sidenote{ The term \emph{Affective Interaction} was explicitly articulated in HCI by \citet{lottridge2011affective}. More recent work has used it to bring together a broader tradition of research that approaches affect as embodied, situated, and socially mediated, in contrast to approaches
centred primarily on emotion recognition and classification \cite{ahmadpour2025affective}.}  Applied to HBI, this lens brings into view how indoor conditions acquire significance through bodily sensation, activity, spatial position, social circumstances, anticipation, prior experience, and occupants' interpretations of technological behaviour.

The analysis identifies five Ways through which occupants encounter and interpret interactive indoor environments: \textit{Interpreting the Somatic Experience Over Time}, \textit{Anticipating Disruption Before It Happens}, \textit{Constructing Spatial Purpose}, \textit{Micro-Adjustments and Low-Effort Escalation}, and \textit{Making Sense of Smartness} (cf.~\autoref{sec:ch3-higher-order}). Together, the Ways show that occupants relate to indoor environments prospectively and interpretively, rather than responding only to conditions as they occur

These findings are consolidated into six Affective Practices, each paired with an enabling Interaction Quality:
\textit{AP1---Orient and Appraise} with
\textit{IQ1---Somatic Legibility};
\textit{AP2---Select a Micro-Condition} with
\textit{IQ2---Spatial Niche Discovery};
\textit{AP3---Social Attunement and Coordination} with
\textit{IQ3---Social Cue Surfacing};
\textit{AP4---Modulate and (De)Escalate} with
\textit{IQ4---(De)Escalation-Aware Interaction};
\textit{AP5---Probe and Align with Smart Building Technology} with
\textit{IQ5---Smartness Legibility}; and
\textit{AP6---Settle and Stabilise} with
\textit{IQ6---Active Non-Intervention} (cf.~\autoref{sec:ch3-objectives}). 

The Affective Practices describe recurring ways in which occupants manage their relationship with indoor conditions, while the Interaction Qualities translate these findings into capacities that interactive environments might support.

The second dissertation contribution therefore extends Affective Interaction (AfI) into HBI as a methodological and conceptual lens for investigating situated indoor experience. It shows how affect can provide an entry point into how environmental conditions come to matter for occupants, while the Affective Practices and Interaction Qualities carry this understanding towards design. The conceptual and methodological commitments of this contribution are developed further in \autoref{sec:affect-contribution}.

\clearpage

\paragraph{\rqref{biophilic}.}

\emph{What possibilities does multisensory biophilic design offer for future interactive indoor environments, and how can these possibilities be articulated to support HBI design?}

The architecture-facing studies presented in \autoref{ch:study-three} and \autoref{ch:study-four} identify multisensory biophilic design as a design direction for extending interactive indoor environments. First, the expert interviews~\autoref{ch:study-three} expand the biophilic design space through eight themes: \textit{Ecological Authenticity and the Essence of Nature in Building Design}, \textit{Art as an Expression for Biophilic Experiences}, \textit{Biophilic Design for Ecological Empathy}, \textit{Multisensory and Multidimensional Biophilic Design}, \textit{Multispecies Flourishing in Buildings}, \textit{Biophilic Design for Social Justice}, \textit{The Body in Biophilic Design}, and \textit{Biophilic Design for Occupant (Dis)Comfort}. These themes expand the experiential and ecological vocabulary available for conceiving biophilic interaction in HBI.

Five design opportunities are also developed with supporting emerging technologies: \textit{Designing for Emergent Natural Processes}, \textit{Designing for Occupant Body Sensitivity}, \textit{Designing for Biophilic (Dis)Comfort}, \textit{Designing for Biophilic Equity}, and \textit{Designing for Relational Encounters with Nature}. These opportunities demonstrate how interactive indoor systems and specific technologies might support novel biophilic experiences.

The subsequent expert design sessions presented in \autoref{ch:study-four} address how such possibilities can be structured for design. Their analysis produces a three-layer Biophilic Interaction Design Framework (BIDF) comprising seven design dimensions, seven design values, and six interaction forms (cf.~\autoref{sec:ch6-overview}). The dimensions structure the range of considerations that designers and researchers can attend to when developing interactive biophilic environments. These are: \textit{Sensory Modality}, \textit{Spatial Context}, \textit{Interaction Type}, \textit{Affective Aim}, \textit{Natural Reference}, \textit{Temporal Behaviour}, and \textit{Level of Mediation}. 

The design(er) values capture the considerations that inform how these possibilities are judged and prioritised i.e., why particular configurations matter. The seven values are: \textit{Sensory Ecology}, \textit{Attentional Ease}, \textit{Proprioceptive Engagement}, \textit{Irreversibility}, \textit{Unknowability}, \textit{Non-Reciprocity}, and \textit{Collective Liability}.
 
The interaction forms describe recurring logics through which occupants might encounter and participate in biophilic indoor environments. These are: \textit{Anchoring}, \textit{Synchronising}, \textit{Dwelling}, \textit{Inferring}, \textit{Withholding}, and \textit{Accepting}. Together, the three layers make expert architectural reasoning available
to HBI---HCI as a structured design resource. In particular, the framework extends consideration beyond immediate responsiveness and personalisation towards interaction that may involve duration, ambiguity, limited control, environmental variation, and shared consequence.

The third dissertation contribution therefore combines an expanded
account of what multisensory biophilic interaction might encompass with
a framework for reasoning about how such interactions might be designed
at architectural scale.
\clearpage
\paragraph{\rqref{sound}.}

\emph{How can an AI-transformed office soundscape be developed as a form
of biophilic interaction, and what does an affective lens reveal about
how it is experienced and interpreted in shared workplace contexts?}

TThe sound-focused direction developed in the third expert design session
in \autoref{ch:study-four} considered how ongoing office sound could be
transformed into nature-like sound while preserving a perceptible
connection to the surrounding office soundscape. This direction informed
the development of the Mediated Office Soundscapes System (MOSS) \autoref{ch:study-six}. MOSS uses a neural audio models
trained on nature sounds to transform incoming office audio into bird-like or water-like sound. 

This relationship between the transformed sound and ongoing office activity is conceptualised as \emph{biophilic mediation}. Biophilic interaction is therefore produced through the ongoing activity of the workplace itself. MOSS provides a working system through which this proposition can be experienced and examined.

The system was examined through an online listening study and live
demonstrations with architecture and spatial-sound experts. The
listening study showed that the AI-transformed soundscape could support
calmness and reduce the clinical character and distraction associated
with the office soundscapes. Participants' accounts further
showed that acceptance depended on how the transformed sound related to
the office context, how strongly it occupied attention, and whether its
response to surrounding activity could be understood.

An affective lens therefore shows that the experiential value of the
soundscape depends on more than the acoustic quality of the generated
nature-like sound. A transformation can remain recognisably artificial
while still fitting the workplace and altering its affective character.
The findings identify three conditions for biophilic mediation in shared
workplaces. First, \emph{contextual coherence}: listeners need to perceive
a meaningful relation between the transformed sound, the surrounding
office activity, and the setting in which it occurs. Second,
\emph{peripheral presence}: the transformation needs to remain
perceptible without continually demanding attention. Third,
\emph{temporal and interpretive intelligibility}: listeners need
sufficient continuity to form an understanding of why and when the
soundscape changes. Artefacts, acoustic clutter, excessive prominence,
or difficult-to-read changes can undermine these conditions and make
the transformation distracting or confusing. The findings further
indicate that such responsiveness should remain configurable,
particularly in shared environments, so that occupants can influence
properties such as prominence, auditory character, and degree of
transformation.

The fourth dissertation contribution therefore comprises MOSS as a
working instance of biophilic mediation and an account of the
experiential conditions that shape how such mediation is encountered in
offices. It shows how ongoing indoor sound can become
material for biophilic interaction and, through an affective lens, how
the success of that interaction depends on contextual fit, attentional
presence, intelligible change, and occupants' ability to negotiate the
soundscape in use.

\paragraph{Responding to the overarching research question.}

This dissertation shows that an affective orientation enables HBI to
specify and investigate experiential complexity by examining how indoor
conditions acquire significance within situated experience. It directs
attention beyond specific environmental variables towards how
occupants sense and interpret surrounding environmental conditions, anticipate change, negotiate available possibilities, and adjust their relationship with their environment over time.

Across the dissertation, this orientation connects three forms of work:
mapping the experiential concerns already addressed within HBI,
investigating how those concerns are constituted in situated use, and
carrying this knowledge into the development and examination of future
interactive environments. The architecture-facing studies extend this
last step by broadening and structuring possible forms of biophilic
interaction, while MOSS provides one technically realised instance
through which such an interaction can be experienced and examined.
Affective HBI therefore connects the investigation of situated indoor
experience with design reasoning about future interactive environments.



\section{affective hbi as an orientation for situated inquiry}
\label{sec:affect-contribution}
This section develops the inquiry-oriented contribution of affective
HBI. Drawing primarily on the occupant studies in
\autoref{ch:study-two}, it considers emotion as an entry point
into situated indoor experience, Affective Practices as a unit of
analysis, and situated methods for retaining the circumstances through
which experience acquires meaning. It then considers how experiential
findings can be carried forward as questions for design.

Affective HBI builds on HBI's concern with architectural-scale interaction
\cite{alavi2019introduction, alavi2018artifacts, alavi2017comfort,
wiberg2008environment, wiberg2020interaction} and on experience-centred HCI's attention to embodied, situated, and felt experience
\cite{dourish2001action,mccarthy2004technology,hook2018designing}.

This position draws particularly on Affective Interaction, which understands emotion as situated, embodied, and constituted through interaction rather than as discrete internal states to be identified independently of context \cite{boehner2005affect,boehner2007emotion,lottridge2011affective, hook2009affective,ahmadpour2025affective}.  Affective HBI extends this
perspective to architectural-scale experience by using affect as an
entry point for investigating how indoor conditions come to matter for
occupants (cf.~\autoref{sec:ch3-s1-findings}, \autoref{sec:ch3-s2-findings}, and \autoref{sec:ch3-disc-reflections-afi}).


Within this dissertation, affective HBI is specified through three
elements of inquiry. Its \emph{object of inquiry} is the situated relation through which an indoor condition acquires affective significance (cf.~\autoref{sec:ch3-s1-findings} and \autoref{sec:ch3-s2-findings}). Its \emph{unit of analysis} is the practice through which an occupant appraises, anticipates, interprets, adjusts to, or settles within that condition (cf.~\autoref{sec:ch3-higher-order}). Its \emph{methodological approach} is to retain the spatial, bodily, technological, social, and temporal circumstances through which experience is produced and interpreted (cf.~\autoref{sec:ch3-s1}, \autoref{sec:ch3-study2}, and \autoref{sec:ch3-objectives}). These three elements are connected to design through a further step: translating experiential relations into questions about what future interactive environments might need to make perceptible, understandable, negotiable, adjustable, or possible. Together, these elements specify how affective HBI produces situated
knowledge about indoor experience and how that knowledge can be carried
forward towards design.


\subsection{emotion as an entry point for investigating indoor experience}
\label{subsec:emotion-entry-point}

Within the occupant studies, emotion is used as a way of investigating
how environmental conditions come to matter for occupants, rather than
as a state to be classified (e.g commonly used in ~\cite{picard1999affective, gao2020n}). Reports of feeling calm, exposed, frustrated, uncertain, or relieved can direct attention to the circumstances that produced those feelings and to how occupants interpreted them (cf.~\autoref{sec:ch3-s1-findings} and \autoref{sec:ch3-s2-findings}). The in-situ survey from \autoref{ch:study-two} supports this methodological role: emotion-framed accounts contained richer bodily, social, and situational information than comfort-framed accounts (cf.~\autoref{sec:ch3-s1-findings}).

The meaning of an emotion depends on the situation in which it is
experienced. Calmness, for example, may relate to quietness, daylight,
enclosure, freedom from interruption, or confidence that conditions
will remain stable, and these qualities matter differently when an
occupant is concentrating, resting, meeting someone, or moving through
the building. Likewise, quietness may support concentration in one
situation while making another person's movements more conspicuous in
another. Affect therefore cannot be assigned a fixed meaning to a
particular environmental condition \cite{feldman2022affect,barrett2022context,boehner2007emotion}.

Affective HBI uses emotion-framed inquiry to examine these relations.
This complements approaches that measure or model affective responses
to environmental stimuli \cite{mehrabian1974approach,russell1980circumplex,
bower2019impact,picard1999affective,gao2020n} by asking how a condition becomes significant within situated use
\cite{boehner2005affect,boehner2007emotion, lottridge2011affective} (cf.~\autoref{sec:ch3-afi} and \autoref{sec:ch3-disc-reflections-afi}).


\subsection{affective practices as the unit of analysis}
\label{subsec:affective-practices}

Emotion provides an entry point into situated experience, while
Affective Practices provide the unit of analysis. They describe the
practical and interpretive work through which occupants establish,
maintain, or revise their relationship with an indoor environment.
This includes attending to bodily and environmental cues, anticipating
change, interpreting spatial and social conditions, testing possible
adjustments, and deciding when a situation is workable
(cf.~\autoref{sec:ch3-s2-findings} and
\autoref{sec:ch3-higher-order}).

This practice-based account shifts analytical attention from relying on objective, reported affective state (e.g. rating valence versus arousal \cite{raodriving}) to how occupants act and make judgements over time. Previous encounters shape expectations, anticipated activity influences
present spatial decisions, and small adjustments allow occupants to test
whether a condition supports what they are trying to do. The
corresponding Interaction Qualities translate these observed practices
into characteristics that interactive environments might support
(cf.~\autoref{sec:ch3-higher-order}).

Affective Practices also place comfort within a wider field of situated
judgement. Environmental conditions and reports of satisfaction remain
important for healthy, safe, sustainable, and usable buildings
\cite{alavi2017comfort,jazizadeh2014human,marques2020real,
van2010effect,zhang2012factors}, but occupants judge these conditions
alongside privacy, activity, social proximity, familiarity, available
alternatives, and expectations of building behaviour. A thermally
imperfect location, for example, may still be preferred because it
provides greater visual protection. The contribution is therefore to
show how comfort forms one part of the broader relations through which
occupants determine whether an environment works for them
(cf.~\autoref{sec:ch2-findings-experiences},
\autoref{sec:ch2-implications-landscape}, and
\autoref{sec:ch3-s2-findings}).


\subsection{situated methods for architectural-scale experience}
\label{subsec:situated-methods}

Architectural-scale experience extends across rooms, routes, spatial arrangements, technological infrastructures, and encounters over time \cite{alavi2019introduction,alavi2018artifacts, wiberg2008environment, wiberg2020interaction,becerik2022field}. Affective HBI therefore requires methods that retain the circumstances in which experience occurs rather than separating reported feelings from situations in which they acquire meaning \cite{dourish2001action,mccarthy2004technology, lottridge2011affective,ahmadpour2025affective}.

The occupant studies combined emotion-framed surveys, building walks, body mapping, situated discussion, and open-ended accounts \cite{mitchell2020no,crivellaro2015contesting, turmo2023towards,hook2018designing} (cf.~\autoref{sec:ch3-s1} and \autoref{sec:ch3-study2}). These methods allowed participants to locate experiences within particular parts of the building, compare spatial alternatives, demonstrate how they responded to conditions, recall previous encounters, and anticipate what might happen next. They therefore connected affective accounts to the situations in which those accounts were formed (cf.~\autoref{sec:ch3-s1-findings} and \autoref{sec:ch3 s2-findings}).

This is particularly important in interactive indoor spaces, where
occupants often encounter technology indirectly through changes in
lighting, sound, temperature, access, or other environmental conditions,
without necessarily knowing the sensing and control processes behind
them \cite{saputra2023smart,daissaoui2020iot} (cf.~\autoref{sec:ch3-s2-findings} and \autoref{sec:ch3-higher-order}). Situated methods make it possible to examine how occupants interpret such behaviour and how those
interpretations influence subsequent action \cite{dourish2001action, bellotti2002making, hu2023scoping,lim2009and}. Their methodological contribution to HBI is therefore to provide explanatory depth while retaining the architectural circumstances through which experience acquires meaning (cf.~\autoref{sec:ch3-objectives}, \autoref{sec:ch3 discussion}, and \autoref{sec:ch3-disc-reflections-afi}).


\subsection{from experiential findings to design questions}
\label{subsec:experiential-design-knowledge}

The preceding elements specify how affective HBI investigates situated
indoor experience. Their relevance to design lies in carrying the
relations identified through that inquiry forward as design questions.
Affective inquiry does not determine a corresponding design response.
An account of feeling exposed, for example, may reveal a relation among
visibility, activity, social presence, and available spatial
alternatives. The resulting design question concerns how those
relations might be supported, altered, made legible, or deliberately
left open, rather than how the reported emotion itself should be
removed.

The Affective Practices and Interaction Qualities provide one
illustration of this translation
(cf.~\autoref{sec:ch3-higher-order}). Practices describe how occupants
orient, interpret, negotiate, adjust, and settle within environmental
conditions; the corresponding Interaction Qualities identify
capacities that an interactive environment might support. They
therefore preserve the relation identified through inquiry while
shifting the question from what an occupant experienced to what a
future environment might enable.

\clearpage
\section{architecture as design knowledge for affective hbi}
\label{sec:architecture-design-knowledge}

Architectural knowledge provides the next step in developing possible
responses to the design questions produced through affective inquiry. HBI has called for interaction design to engage buildings as persistent, shared, and spatially distributed settings of computation \cite{alavi2018artifacts, alavi2019introduction, wiberg2020interaction} and to work with the practitioners and domain experts who shape them \cite{kirsh2019architects}. The dissertation takes up this direction through sustained engagement with architectural and spatial expertise across expert interviews (cf.~\autoref{sec:ch4-methodology}), expert design sessions (cf.~\autoref{subsec:ch6-study-design}), and live prototype demonstrations (cf.~\autoref{ch:study-six}).

Architecture contributes here not just as the setting in which interaction takes place, but as a source of design knowledge. Architectural expertise brings ways of reasoning about how spatial, sensory, material, environmental, and temporal conditions are composed and inhabited \cite{pallasmaa2024eyes, griffero2016atmospheres, zumthor2006atmospheres, verma2022appraisal}. The architecture-facing studies bring this knowledge into conversation with embodied, experience-centred \cite{hook2018designing}, soma-based \cite{tschaepe2021somaesthetics}, and material and spatial approaches in HCI \cite{van2018skin,nabil2017interioractive}. This exchange broadens the forms of knowledge through which HBI can conceive and judge future interactive environments. 

The section develops this contribution in two steps. It first examines how architectural expertise can function as a design resource within HBI, and then considers how this exchange changes the object of interaction design at architectural scale.


\subsection{architectural expertise as a design resource}
\label{subsec:architectural-expertise}

Engaging architectural expertise provides access to forms of domain knowledge developed through architectural practice and creates ways of bringing that knowledge into interaction-design inquiry. In the expert interviews (cf.~\autoref{sec:ch4-methodology}), architectural and spatial practitioners articulated knowledge, experience, and future possibilities that expanded the design questions available to HBI. In the subsequent design sessions (cf.~\autoref{subsec:ch6-study-design}), HCI design methods provided a means of externalising and examining this knowledge through concrete design proposals.

The Biophilic Interaction Design Framework (cf.~\autoref{sec:ch6-overview}) is one outcome of this exchange. Its dissertation level significance lies in demonstrating that forms of architectural practitioner knowledge can be elicited and structured as a design resource for HCI. Rather than treating architects only as informants about the built context, the process makes their design judgements available for analysis, comparison, and further development within interaction design.

The subsequent live demonstrations of MOSS (cf.~\autoref{ch:study-six}) extended this exchange to an emerging form of AI-based interaction. They asked architectural and spatial-sound experts to consider computational behaviour itself as a matter of spatial design, rather than treating the technology as something simply accommodated within a finished architectural setting. This raised questions of where such behaviour should be located, how prominently it should be perceived, how it should relate to ongoing activity, and how occupants might understand or influence it. \sidenote{This boundary is already beginning to shift in architectural practice. In discussions surrounding the studies, some practitioners described recent projects for AI technology-based companies in which interactive lighting and other responsive systems formed part of workplace design. Such examples suggest increasing architectural engagement with computationally mediated environments, although these forms of interaction are not yet routine across practice.}

Across these engagements, knowledge moved in both directions. Architectural practitioners brought domain knowledge from architectural practice into HCI research, while HCI methods provided ways of eliciting, structuring, technically developing, and critically examining that knowledge.  The collaboration therefore concerns more than applying interactive technologies within buildings, it provides a means through which architectural and interaction-design knowledge can shape one another. This offers one response to HBI's call for closer engagement with the disciplines responsible for the design of built environments.

\subsection{architectural-scale interaction as an inhabitable relation}
\label{subsec:architectural-inhabitable-relation}

The reciprocal exchange described above also changes what interaction design treats as its object at building scale. Architectural expertise directs attention beyond the behaviour of a computational system towards the environmental conditions produced through its operation. Computational behaviour is encountered as part of an environment that occupants inhabit alongside spatial conditions, environmental processes, other people, and ongoing activities \cite{kirsh2019architects,alavi2019introduction}. I use the term \emph{inhabitable relation} to describe this broader unit of interaction.

An inhabitable relation is spatially distributed and temporally extended \cite{alavi2018artifacts,nguyen2024adaptive}. A computational change may alter conditions beyond the person who initiated it, remain perceptible after the initiating action has passed, or acquire different significance as occupants move, activities change, and other people enter the setting. Studies of shared environmental systems likewise show that interaction with environmental controls becomes a matter of coordination among occupants rather than solely an individual interaction \cite{niemantsverdriet2018share}. Interaction at building scale therefore has to be considered in relation to the persistent and shared character of the environment \cite{alavi2019introduction, wiberg2020interaction}, rather than only through the correspondence between an input and a system response.

This changes the scope of interaction-design judgement. The relevant question is not only how well a computational system responds, but what kind of environmental condition its behaviour produces and how that condition can be lived with over time. Research on responsive architecture similarly argues that adaptation should be considered in relation to the spatial, situational, and subjective qualities of the place it produces \cite{nguyen2022responsive}. Timing and spatial distribution \cite{nguyen2024adaptive}, as well as consequences for other occupants \cite{niemantsverdriet2018share}, therefore become matters of interaction design alongside the behaviour of the system itself. The architectural perspective extends HBI from designing computational behaviour towards designing the conditions through which that behaviour becomes part of an occupied environment.

This distinction is particularly consequential for responsive environments. Moving beyond conventional interfaces can make questions of what a system is doing and how its actions should be interpreted newly salient \cite{bellotti2002making}, while shared environments require automatic behaviour to coexist with multiple occupants and forms of control \cite{vanDeWerff2019lighting}. A response may therefore be technically appropriate while remaining difficult to interpret, intrusive within an activity, disruptive to others, or poorly integrated with the surrounding environment. Responsiveness at architectural scale must therefore be considered in relation to the conditions of inhabitation in which computational change is encountered. The following section examines one such relation through the technically realised case of MOSS.



\section{moss as a design instance}
\label{sec:sound-instance}

The preceding section established architectural knowledge as a resource for developing interactive environments at building scale. This section examines one sound-focused biophilic direction developed through the expert design sessions (cf.~\autoref{subsec:ch5-session3}) as a technically realised design instance through MOSS. It provides a concrete case through which an environmental interaction can be experienced and examined.

The affective orientation \cite{boehner2005affect,boehner2007emotion} developed in \autoref{sec:affect-contribution} shapes this examination. Attention is directed towards how listeners relate the transformed sound to surrounding activity, how system behaviour enters attention, and whether the resulting condition fits the activity and workplace situation (cf.~\autoref{ch:study-six}).  MOSS thereby brings the dissertation's biophilic design work and affective inquiry together in a technically realised system.


\subsection{biophilic mediation as an interaction design concept}
\label{subsec:biophilic-mediation}

MOSS develops \emph{biophilic mediation} as an interaction design
concept. I use \emph{mediation} to describe an altered relationship
between a listener and an existing environmental condition. In MOSS,
sounds occurring in the indoor environment are transformed into
nature-like sound while characteristics of the incoming audio continue
to shape what the listener hears (cf.~\autoref{ch:study-six}). The system therefore changes how ongoing acoustic activity is encountered while retaining a perceptible, variable, and partly ambiguous connection to that activity.

This distinguishes biophilic mediation from conventional nature
playback, where an independent nature recording is introduced into the
soundscape, and from acoustic masking, which commonly seeks to reduce
the prominence of unwanted sound \cite{haas2018soundscape, jacobsen2024personal}.

Biophilic mediation was developed around three \emph{design aims}:
\emph{enriching quietness} through low-level nature-like variation,
\emph{reframing activity} by changing the auditory character of office
events, and \emph{preserving environmental connection} by allowing
changes in surrounding activity to remain consequential to the
soundscape
(cf.~\autoref{ch:study-six}). These aims define the intended form of
interaction. The following subsection identifies the experiential
conditions that shaped how listeners encountered it.


\subsection{experiential conditions of biophilic mediation}
\label{subsec:moss-experiential-conditions}

The studies identify three connected experiential conditions that shape how biophilic mediation is encountered in an office context: \emph{contextual integration}, \emph{bounded salience}, and \emph{temporal and interpretive coherence} (cf.~\autoref{ch:study-six}). These conditions connect the preceding design aims to how listeners experienced and interpreted the realised soundscape.

The first condition is \emph{contextual integration}. In the listening study, the AI-transformed sound was rated as less nature-like than conventional nature sound while receiving comparable pleasantness ratings, higher calmness ratings, and lower ratings for cold or clinical character, distraction, and being out of place (cf.~\autoref{ch:study-six}). Participants also recognised the intended natural referents while reporting artificiality, artefacts, restricted variation, and discontinuity (cf.~\autoref{ch:study-six}).

These findings distinguish natural realism from contextual integration as two qualities of the listening experience. A recognisably artificial transformation can establish a coherent relationship with office activity, while an independent nature recording can remain experientially detached from that activity. \emph{Contextual integration} concerns how the transformed sound relates to the spatial, temporal, social, and activity conditions in which it is encountered. For biophilic mediation, this relationship with the surrounding environment constitutes an experiential quality alongside the recognisability and realism of the natural reference.

The second condition is \emph{bounded salience}: the transformed sound remains perceptible while occupying a limited position in focal attention. This concern relates to work on calm and unremarkable interaction \cite{case2015calm,tolmie2002unremarkable}. In the listening study, soft, continuous, and proportionate output could alter the felt character of the office while remaining peripheral, whereas sharp events, abrupt transitions, repetition, and excessive sensitivity increased attentional demand (cf.~\autoref{ch:study-six}).

Within the dissertation, bounded salience provides a more specific account of the \emph{Attentional Ease} developed in the Biophilic Interaction Design Framework (cf.~\autoref{sec:ch5-layer2}). It concerns how strongly and how frequently computational behaviour becomes perceptually prominent over time. The soundscape can contribute atmospheric variation and retain a perceptible relationship with activity while allowing occupants' attention to remain directed towards the activities for which the setting is being used.

The third condition is \emph{temporal and interpretive coherence}. Participants differed in whether they perceived the system as responsive and in the relationships they inferred between office activity and transformed output (cf.~\autoref{ch:study-six}). The experience of responsiveness depends on whether listeners can relate changes in the soundscape to preceding events and form a workable account of the system's broad behaviour \cite{bellotti2002making,lim2009and}. In MOSS, coherent relationships between activity and output could support liveliness and environmental connection, whereas highly salient, delayed, inconsistent, or difficult-to-anticipate changes could amplify activity, expose the mechanism, or increase self-consciousness (cf.~\autoref{ch:study-six}). \emph{Temporal and interpretive coherence} therefore concerns whether computational change remains sufficiently connected to ongoing activity for listeners to interpret its behaviour over time.

The expert demonstrations place these experiential conditions in relation to spatial implementation (cf.~\autoref{ch:study-six}). Experts considered personal, localised, room-level, and building-integrated forms of responsive sound and related their suitability to room purpose, acoustic character, activity, duration, and delivery format (cf.~\autoref{ch:study-six}). Their requests for influence over sensitivity, timing, auditory character, intensity, and system operation further show how these conditions are shaped by configuration in shared environments (cf.~\autoref{ch:study-six}). At architectural scale, their realisation therefore depends on where the intervention is encountered, how its behaviour is distributed through space and time, and how occupants can participate in shaping that behaviour.

Together, the three conditions locate the experiential contribution of biophilic mediation in the relationships among computational behaviour, ongoing environmental activity, attention, interpretation, and use. They specify how an AI-transformed soundscape becomes an experiential condition of the workplace and provide the basis for examining how the technical properties of MOSS shape that condition.


\subsection{technical properties as environmental qualities}
\label{subsec:technical-experiential}

MOSS shows how technical properties and implementation decisions become
environmental qualities when computational behaviour forms part of an
inhabited condition \cite{alavi2018artifacts, alavi2019introduction, wiberg2020interaction}. Latency, smoothing, model behaviour, sensing configuration, mixing, and playback configuration shape perceptible qualities of the soundscape, including its continuity, prominence, timing, and the extent to which ongoing office activity remains audible (cf.~\autoref{ch:study-six}).


In MOSS, latency affects whether an auditory change appears related to an office event; smoothing shapes whether the output is experienced as a continuous texture or as a sequence of discrete reactions; and model artefacts can increase attentional demand and alter the perceived naturalness of the sound (cf.~\autoref{ch:study-six}). Microphone placement determines which activities become consequential to the transformation, while playback level and mixing shape the prominence of the transformed sound and the degree to which ongoing office activity remains perceptible (cf.~\autoref{ch:study-six}). These technical decisions therefore participate directly in producing the experiential conditions described in the preceding subsection.

This relationship changes the scope of technical evaluation in HBI. For MOSS, the relevant system includes the microphone input, preprocessing, neural model, temporal conditioning, mixing, playback, room acoustics, surrounding activity, and listening situation \cite{alavi2018artifacts, alavi2019introduction, wiberg2020interaction} (cf.~\autoref{ch:study-six}). Experiential quality is produced through relations among these components. Technical evaluation at building scale therefore includes the sensory, spatial, temporal, and social conditions produced through system behaviour.

The contribution here is specific to the sound-based case examined through MOSS. Other forms of environmental interaction, including lighting, airflow, temperature, projection, and kinetic systems (cf.~\autoref{sec:ch2-findings-interventions}), involve different perceptual, safety, spatial, and infrastructural considerations. The broader implication for environmental AI is that computational properties acquire significance through the environmental conditions they produce and through how occupants perceive, interpret, and act within those conditions.

\subsection{connecting biophilic design and affective inquiry}
\label{subsec:moss-integration}

MOSS brings the dissertation's biophilic design work and affective inquiry together in a technically realised environmental interaction. The biophilic studies inform the form of interaction being developed, including its concern with sensory variation, environmental connection, and the temporal character of indoor experience (cf.~\autoref{ch:study-four} and \autoref{ch:study-five}). The affective orientation informs how that interaction is examined by directing attention towards how listeners perceive and interpret system behaviour, relate it to ongoing activity, and judge its fit within the workplace situation (cf.~\autoref{sec:affect-contribution} and \autoref{ch:study-six}).

This integration shows how affective HBI can carry a design direction from architectural and interaction-design reasoning into technical development and experiential examination. In MOSS, the behaviour of the system becomes part of an environmental condition that occupants must interpret in relation to their activity and surroundings. Their experience therefore depends partly on how computational change becomes perceptible, intelligible, and open to adjustment. This relation between environmental behaviour and occupant interpretation provides the basis for the following section, which develops a participative account of environmental responsiveness.


% \section{a participative account of environmental responsiveness}
% \label{sec:participative-responsiveness}

% Responsive and adaptive environments have long examined how buildings can alter spatial or environmental conditions in relation to occupants, activities, and sensed data. Research has studied environmental automation, adaptive architecture, personal data, kinetic elements, movable partitions, shared control, and occupant-centred building management\cite{schnadelbach2019adaptive,nguyen2022responsive,nguyen2024adaptive}. These systems commonly raise questions of intelligibility, control, privacy, and fit among several users.

% The synthesis across this dissertation qualifies what responsiveness means experientially. A response is not exhausted by the successful detection of an input and production of an output. Its experiential significance depends on how occupants perceive the change, relate it to preceding conditions, interpret its purpose, adjust their conduct, and understand the alternatives available afterwards.

% This conclusion is supported through convergences across the portfolio. The occupant studies describe the interpretive and practical work through which people inhabit existing environments. The architecture-facing studies describe how experts reasoned about future environmental relations. MOSS provides a concrete system in which the connection among activity, computational transformation, and perceptible output can be examined. Across these different forms of evidence, three concerns recur: bodily and interpretive legibility, temporality, and occupant participation.


% \subsection{bodily and interpretive legibility}
% \label{subsec:bodily-interpretive-legibility}

% Occupants frequently assessed indoor environments through bodily and sensory impressions. Light, sound, temperature, enclosure, visibility, furniture, and social density contributed to whether a place felt suitable for a particular activity. These impressions could precede an explicit explanation. A person might feel exposed, settled, constrained, or restless before identifying the combination of conditions producing that judgement.

% The Affective Practice of \emph{Orient and Appraise} and the Interaction Quality of \emph{Somatic Legibility} articulate this relation. The body participates in understanding the environment and in deciding whether to remain, move, or alter an activity.

% The architecture-facing studies give sensory and bodily relations a corresponding design role. Sensory Ecology and Proprioceptive Engagement treat sensory composition and bodily movement as sources of environmental understanding. Anchoring and Inferring describe interaction through contact, posture, movement, gradients, and indirect sensory information.

% These findings connect with embodied and soma-based HCI, which positions felt bodily experience as a site of interaction. Their architectural-scale implication concerns how computational behaviour becomes perceptible through environmental conditions. A building can communicate through changes in sound, airflow, illumination, temperature, spatial configuration, or other sensory effects. Legibility therefore includes the ability to sense that a relevant difference has occurred and to relate that difference to the activity or condition that produced it.

% Complete explanation is neither possible nor desirable for every ambient change. Work on ambiguity has shown how partial interpretation can support reflection, curiosity, and appropriation\cite{gaver2003ambiguity}. The framework's Unknowability value similarly protects variation and partial understanding. The dissertation suggests that environmental interaction requires \emph{calibrated legibility}: sufficient perceptibility for occupants to establish a workable account of consequential behaviour, with room for ambiguity where its consequences remain limited and its experiential purpose is coherent.


% \subsection{temporality and response as relation}
% \label{subsec:response-relation}

% Occupants interpreted buildings through current conditions, prior encounters, and expected changes. They anticipated sunlight, noise, occupancy, and automated behaviour. They tested their expectations by moving, adjusting, waiting, and observing. Through repeated encounters, they developed practical accounts of what a building could do and how reliably it would do it.

% The biophilic framework also locates interaction across time. Synchronising connects activity to environmental rhythms. Dwelling gives sustained presence a role in how experience accumulates. Inferring requires observation across repeated cues. Irreversibility gives significance to traces and consequences that persist.

% MOSS makes the temporal character of responsiveness technically explicit. Latency, continuity, smoothing, persistence, and sensitivity affect whether transformed sound appears connected to office activity, overly reactive, fragmented, or independent from it. The listener interprets responsiveness through a temporal relation between an event and the changing auditory output.

% Responsive interaction should therefore be assessed over a sequence rather than at a single input--output moment. Relevant questions include whether occupants can anticipate the broad character of a response, recognise continuity across events, revise an earlier interpretation, and establish familiarity without the system becoming rigid or repetitive.

% This account extends research on calm and unremarkable computing\cite{case2015calm,tolmie2002unremarkable}. Peripheral systems must remain available to perception while respecting the primary activity underway. Their quality depends on timing, salience, continuity, and the occupant's capacity to incorporate them into routine conduct.


% \subsection{adaptation includes occupant participation}
% \label{subsec:adaptation-participation}

% Adaptation describes how a system changes environmental conditions in response to sensed information. The dissertation shows that occupants also participate in how an adapted condition becomes meaningful and workable.

% In the occupant studies, participation took the form of spatial selection, bodily adjustment, social coordination, direct control, interpretation, anticipation, waiting, acceptance, and avoidance. A person who moves to another acoustic zone, waits to see whether ventilation changes, adjusts a blind, or accepts a condition after comparing alternatives participates in the environmental arrangement.

% The biophilic interaction forms broaden this account. Participation may occur through contact and stability, movement through spatial gradients, sustained presence, attention to indirect cues, deliberate distance, or acceptance of conditions that remain partly outside individual control. These forms show how interaction can involve degrees of involvement beyond command-based control.

% Research on adaptive architecture and shared spatial control has similarly shown that occupants negotiate among system initiative, direct operation, and the competing demands of shared environments\cite{nguyen2022responsive,nguyen2024adaptive}. The dissertation contributes an affective account of this negotiation: participation includes the interpretive and practical processes through which occupants decide whether a response supports the activity and situation at hand.

% Interactive environments can support such participation through perceptible cues, meaningful alternatives, reversible adjustments, spatial variation, and opportunities to observe environmental behaviour before committing to a change. Immediate correction remains important for health, safety, access, and acute discomfort. Other experiential interventions may benefit from slower change, suggestive support, or deliberate restraint.


% \subsection{collective configuration in shared environments}
% \label{subsec:collective-configuration}

% Indoor environments are frequently shared. A response that benefits one occupant may distract, expose, constrain, or inconvenience another. The occupant studies show that people already manage these conditions through deference, negotiation, spatial movement, and informal social judgement. The architecture-facing studies introduce values of Collective Liability and Non-Reciprocity, which draw attention to shared consequence and relations that exceed individual preference. The MOSS studies show that sensory tolerance, task demands, desired sound character, and interpretations of responsiveness differ across participants.

% These findings direct attention from isolated control towards collective configuration. Relevant design concerns include:

% \begin{itemize}
%     \item how occupants understand the broad behaviour of a shared system;
%     \item which parameters can be adjusted and by whom;
%     \item how configurations vary across activities and periods of use;
%     \item which spatial alternatives remain available; and
%     \item how occupants can request, revise, pause, or withdraw from an intervention.
% \end{itemize}

% The dissertation does not establish a model of environmental governance. It raises \emph{shared governability} as a future research concern: how authority, objection, responsibility, and revision should be organised when ambient computational systems alter collectively experienced conditions. This concern involves organisational and institutional arrangements alongside interface and system design.


% \subsection{participative responsiveness as a dissertation-level implication}
% \label{subsec:participative-implication}

% The portfolio supports a participative account of environmental responsiveness. Under this account, the experience of responsiveness is formed through the relation between system behaviour and occupant interpretation, adjustment, and participation.

% This proposition follows from four bodies of evidence:

% \begin{enumerate}
%     \item occupant accounts show that people appraise, anticipate, probe, coordinate, and settle in relation to environmental conditions;
%     \item expert interviews expand future design towards interpretation, participation, variation, and collective consequence;
%     \item the framework articulates interaction forms involving bodily movement, duration, indirect cues, distance, and restraint; and
%     \item MOSS shows that the perceived value of computational response depends on its temporal coherence, attentional role, intelligibility, contextual fit, and configurability.
% \end{enumerate}

% Participative responsiveness is therefore an implication of the four dissertation contributions. It is not presented as a separately validated framework or a fifth contribution. It identifies a direction for HBI: environmental intelligence should be judged through the quality of the relations it permits occupants to establish with computational behaviour.


% \section{scope, transferability, and reflexivity}
% \label{sec:scope-transferability}

% The dissertation develops one affective route into the wider complexity identified by the scoping review. Its claims concern how affect can provide access to situated
% significance, how architectural expertise can expand and articulate
% experiential design possibilities, and how one sound-focused biophilic
% direction can be developed as an AI-transformed office soundscape and
% examined through an affective lens.

% The four contributions produce different forms of knowledge. The scoping review provides descriptive and classificatory knowledge about a research field. The occupant studies provide exploratory empirical and conceptual knowledge concerning situated experience. The interviews provide professional accounts and future-oriented imaginaries. The design sessions provide generative design knowledge. The soundscape studies provide comparative listening evidence and professional critique of an early-stage prototype. The dissertation-level synthesis retains these distinctions.


% \subsection{scope of the occupant knowledge}
% \label{subsec:scope-occupants}

% The principal occupant studies centred on LAB42, a recent university building characterised by shared spaces, flexible patterns of use, environmental sensing and automation, and limited direct occupant control. This setting was strategically informative because it concentrates conditions increasingly present in contemporary institutional and workplace environments.\sidenote{The in-situ survey also included an older university building as a comparison, while the detailed situated inquiry centred on LAB42.}

% The setting enabled investigation of spatial flexibility, shared environmental conditions, social visibility, and partially automated behaviour in everyday use. The Affective Practices and Interaction Qualities are most directly relevant to buildings with comparable spatial, technological, and organisational conditions.

% Homes, healthcare environments, schools, transport settings, and highly regulated workplaces may produce different forms of appraisal, authority, adjustment, and settlement. The constructs should therefore be applied as exploratory concepts whose recurrence and variation require comparative examination.

% The methods provide detailed evidence about situated encounters and participants' accounts of repeated use. They offer limited direct evidence about seasonal change, multi-year adaptation, organisational change, or sustained relationships with newly introduced systems.


% \subsection{scope of the architecture-facing knowledge}
% \label{subsec:scope-architecture}

% The architecture-facing studies were designed to access professional reasoning concerning atmosphere, materiality, spatial behaviour, temporality, biophilic design, and technologically mediated environments.

% The interviews identify expert discourse and imaginaries. Their themes articulate possibilities that participants considered consequential; they do not establish occupant needs or universal principles. The framework structures reasoning elicited through speculative design sessions. Its dimensions, values, and interaction forms support design analysis and ideation, while their application in occupied environments requires collaboration with occupants, accessibility specialists, engineers, building operators, maintenance staff, and affected communities.

% Most experts worked within European professional contexts. The resulting knowledge may reproduce assumptions concerning climate, sensory desirability, seasonality, professional authority, and relations with nature. Justice-oriented and more-than-human concerns were represented through expert interpretation. Direct engagement with disability-led, Indigenous, non-Western, community-based, ecological, and maintenance-centred knowledge remains necessary.


% \subsection{scope of the soundscape evidence}
% \label{subsec:scope-sound}

% MOSS is a deliberately bounded design instance. The online study isolates immediate responses to controlled sound conditions. Explicit labels and the nature-oriented framing of the study may have shaped how participants interpreted the stimuli.

% The live demonstrations provide direct encounters with a headphone-based prototype. Participants tested its response to ordinary office sounds and then considered its affective, biophilic, spatial, and interactional implications. This framing supported focused professional critique and directed attention towards the concerns embedded in the research questions.

% The studies do not establish habituation, organisational acceptance, room-level collective experience, long-term technical stability, or the consequences of continuous use. They identify the conditions and tensions that a deployed study would need to examine.


% \subsection{researcher position and analytic accountability}
% \label{subsec:reflexivity}

% The researcher occupied several roles across the programme: reviewer constructing classifications, field researcher eliciting occupant accounts, interviewer and workshop facilitator shaping expert discussion, analyst interpreting verbal and visual material, and designer developing and tuning MOSS.

% These roles influenced what became visible. Emotion-framed prompts oriented occupants towards affective interpretation. Biophilic framing directed expert discussion towards sensory experience, nature, and ecology. The description of MOSS as AI-transformed and nature-like provided necessary study context while shaping possible interpretations.

% Analytic accountability therefore depends on preserving the relationship between claims and their sources. The occupant constructs were developed from situated accounts. The interview themes represent professional discourse and imaginaries. The framework derives from verbal reasoning, design concepts, sketches, and other workshop artefacts whose session context was retained during analysis. The soundscape findings distinguish comparative listening evidence from interpretations and projected applications. Contradictory and critical accounts remain consequential because they identify where the intended proposition failed to achieve coherence or fit.

% The convergences developed in this chapter were constructed through dissertation-level comparative interpretation. The studies were undertaken for different purposes and do not constitute a single comparative experiment. Their connections should be read as argued relations among empirical findings, expert propositions, and design implications.

% A different programme could have pursued other experiences, disciplinary partnerships, design approaches, or technological media.The contribution of this dissertation lies in the traceable route it
% develops from mapping how experience is addressed through design and
% technology, through affective inquiry and architecture-facing design
% research, to the development and examination of a sound-focused
% biophilic system and the broader implications drawn across these
% studies.


% \section{future research}
% \label{sec:future-research}

% The dissertation establishes three connected programmes for further research.


% \subsection{extending affective hbi}
% \label{subsec:future-affective}

% Comparative and longitudinal research should examine which Affective Practices and Interaction Qualities recur across building types and which depend on particular spatial, cultural, technological, and institutional conditions. A central question is:

% \begin{quote}
% Which processes of appraisal, anticipation, interpretation, adjustment, and settlement remain recognisable across settings, and how are they altered by authority, sensory diversity, available alternatives, and duration of occupation?
% \end{quote}

% This work can combine situated qualitative methods with environmental, behavioural, and physiological evidence while retaining the interpretive and contextual role assigned to each form of data.


% \subsection{applying multisensory biophilic design knowledge}
% \label{subsec:future-biophilic}

% The biophilic framework requires application and revision through concrete, occupied design processes. Its dimensions, values, and interaction forms should be examined when speculative propositions encounter accessibility, engineering, regulation, maintenance, climate, organisational practice, and everyday use. A central question is:

% \begin{quote}
% Which parts of the framework remain useful when architects, interaction designers, occupants, operators, and other affected communities collectively translate its design knowledge into built interventions?
% \end{quote}

% Research across sound, light, airflow, temperature, smell, haptics, material movement, and multisensory combinations can clarify how architectural reasoning survives, changes, or disappears during technical development.


% \subsection{studying participative environmental ai}
% \label{subsec:future-participative-ai}

% Participative responsiveness requires examination through deployed systems whose effects are experienced over time and by several occupants. Research should compare arrangements involving direct control, spatial choice, reversible settings, environmental gradients, peripheral cues, collective configuration, and deliberate restraint. A central question is:

% \begin{quote}
% Under what conditions can occupants understand, revise, and contest their relationship with ambient AI systems whose behaviour is continuous and whose consequences are spatially shared?
% \end{quote}

% Long-term MOSS deployments could examine source connection, temporal coherence, artefacts, room acoustics, sensory tolerance, configurability, habituation, privacy, and organisational responsibility. Such deployments would also provide the empirical basis required to develop shared governability beyond the concern raised by this dissertation.


% \section{conclusion}
% \label{sec:final-conclusion}

% Interactive buildings increasingly sense, infer, and act. These capacities acquire significance through the conditions they create for the people who inhabit them. A change in light, sound, airflow, access, or temperature is experienced in relation to the activity underway, the people affected, the expectations formed, and the possibilities available for response.

% Affective HBI places these relations at the centre of inquiry and design. It studies how occupants sense environmental conditions, interpret spatial and computational behaviour, anticipate change, coordinate with others, and establish workable forms of inhabitation. Architectural and multisensory biophilic knowledge extends this orientation towards the design of environments whose experiential qualities are composed across space, bodies, technologies, natural processes, and time.

% The dissertation consequently supports a participative account of environmental responsiveness. The quality of an interactive environment depends on the relation between computational behaviour and the occupants who must understand, accommodate, influence, and live with its consequences. A building becomes meaningfully interactive when its computational capacities help people make sense of, act within, and belong to the places they inhabit.




